% Generated draft: requires mathematical and source review.
\documentclass[11pt]{article}
\usepackage[margin=25mm]{geometry}
\usepackage{fontspec}
\setmainfont{Latin Modern Roman}
\usepackage{amsmath,amssymb,mathtools,unicode-math}
\setmathfont{Latin Modern Math}
\usepackage{enumitem,needspace,xurl,hyperref}
\setlength{\parindent}{0pt}
\setlength{\parskip}{5pt}
\setlength{\emergencystretch}{3em}
\newcommand{\N}{\mathbb{N}}
\newcommand{\Z}{\mathbb{Z}}
\newcommand{\Q}{\mathbb{Q}}
\newcommand{\R}{\mathbb{R}}
\newcommand{\C}{\mathbb{C}}
\newcommand{\F}{\mathbb{F}}
\newcommand{\A}{\mathbb{A}}
\newcommand{\PP}{\mathbb{P}}
\newcommand{\E}{\mathbb{E}}
\newcommand{\eps}{\varepsilon}
\newcommand{\dd}{\,\mathrm{d}}
\begin{document}
\section*{UnsolvedMath ID: 11000158}
\textbf{Source problem code:} catalogue entry
\subsection{Definitions and mathematical statement}
Assume, as the source explicitly requires, that \(M\) is a \(3\)-manifold with nonempty boundary. A conventional compression body \(C\) is formed from \(F\times[0,1]\) by attaching \(1\)-handles to \(F\times\{1\}\), with negative boundary \(\partial_-C=F\times\{0\}\) and positive boundary \(\partial_+C\); handlebody conventions cover the empty-negative-boundary case. The source does not supply the relevant handle decomposition, surface markings or boundary-equivalence category.

The preceding double-coset setup \(H\phi H\) is also missing. Conventionally \(H\) is the subgroup of a Heegaard-surface mapping class group whose classes extend over a handlebody. For a marked compression body, the corresponding extension subgroup would consist of boundary mapping classes extending over \(C\), with the required behavior on \(\partial_-C\). Depending on the gluing convention, replacing one handlebody side gives a double coset of the schematic form
\[
 H_C\phi H\qquad\text{or}\qquad H\phi H_C.
\]
These are convention-dependent descriptions, not a recovered definition of the source's groups.

Problem 2.1 asks to make this modification precise from the compression body's handle decomposition, and then determine what happens for a knot space. It requires identifying the actual subgroup/actions and equivalence relation encoded by the modified coset, not merely naming an extension subgroup. In the usual knot-exterior interpretation, \(S^3\setminus\operatorname{int}\nu K\) has torus boundary, so that boundary must be retained in the compression-body setup; whether meridian/longitude data or boundary markings must be fixed is not stated. The source does not define knot space or fix the ambient manifold. No specific double-coset classification or subgroup generators are supplied.
\subsection{Short English statement}
Adapt the source's handlebody double-coset description to a three-manifold with boundary by replacing one side with the correct compression-body subgroup, using its handle decomposition. Work out the knot-space case. The original setup and peripheral-marking conventions are missing.
\subsection{Sources}
\begin{itemize}
\item Local UnsolvedMath registry, permanent ID \(11000158\), exact source code \texttt{catalogue entry}, Problem 2.1, PDF page 156. No individual proposer or proposed year is supplied.
\item Farb, edited volume \emph{Problems on Mapping Class Groups and Related Topics}: \url{https://www.math.uchicago.edu/~farb/papers/mcgbook.pdf}. The volume has 329 explicitly typeset numbered items across chapters, repeated local numbering distinguished by PDF page.
\item Extraction pdf-font-structure, accessed 2026-07-29. Source evidence NEEDS\_REVIEW presents an open problem and calls for a current primary-literature audit. Registry status open and classification OPEN-TRIAGE; the research summary says only ``Literature status: Truncated.'' No separate reference for the preceding double-coset construction is supplied.
\item Dataset attribution: Ulam AI / UnsolvedMath Contributors, CC BY 4.0. Source rights NEEDS\_REVIEW: publicly accessible author/publisher copy, redistribution terms requiring review.
\end{itemize}
\end{document}
